\scanpage{622}
\mspage{10}
\deleted{\unsure{Soit}} Si $\Psi$ est une partie de $\mathrm{Fl}\,M$, on en déduit
pour \unsure{tte} $I$ dans $\mathrm{Cat}$ une partie

\[
\left\{\begin{aligned}
&\Psi(I\unsure{\rangle|)}\subset\mathrm{Fl}\,M(I),\qquad\text{avec }M(I)=\underline{\mathrm{Hom}}(I^{\mathrm o},M),\\
&f\in\Psi(I)\ \Longleftrightarrow\ f(i)\in\Psi\quad \forall i\in\mathrm{Ob}\,I.
\end{aligned}\right.\tag*{(1.11)}
\]

\deleted{\unsure{Soit}} Si on introduit la catégorie discrète $I_0$
ss-jacente à $\mathrm{Ob}\,I$, et le foncteur canonique

\[
\unsure{v|u}_I=v:I_0\longrightarrow I,\tag*{(1.12)}
\]

on trouve un couple de foncteurs adjoints

\diagram{\begin{tikzcd}[ampersand replacement=\&,row sep=0pt,column sep=large]
\text{(1.1\deleted{2}3)} \& M(I)\arrow[r,shift right=1.2,"v^*"'] \& M(I_0)\simeq M^{I_0}\arrow[l,shift right=1.2,"v_!"',"{\unsure{\cdot|>}}"{description,pos=.08}] ,\\
\& {[}\phantom{M(I)} \& \phantom{M(I_0)}{]}\arrow[l,dotted,"v_*"]
\end{tikzcd}}

et on a

\[
\Psi(I)=(v^*)^{-1}\big(\Psi(I_0)\big)\unsure{;|.}
\]

\deleted{donc par 1.6}
\[
\deleted{\Psi(I)=v_!(\Phi(I_0))}
\]

Supposons maintenant qu'on ait

\[
\left\{\begin{aligned}
&\Psi=\Phi_*\\
&\text{d'où}\qquad\Psi(I_0)=\Phi(I_0)_*
\end{aligned}\right.\tag*{(1.14)}
\]

(vérification immédiate), on conclut alors par 1.6

\[
\Psi(I)=\big(v_!(\Phi(I_0)\big)_*=\big(\Phi\langle I\rangle\big)_*\tag*{(1.15)}
\]

avec

\[
\Phi\langle I\rangle\overset{\mathrm{def}}{=}v_!\big(\Phi(I_0)\big)^{\sim}=\Psi(I)^*.\tag*{(1.16)}
\]

Si on a un foncteur

\[
u:I\longrightarrow I',\tag*{(1.17)}
\]

d'où un triple de foncteurs adjoints,
% archive sheet number: 621
